\documentclass[11pt]{article}
\usepackage{amsmath,amssymb,amsthm}
\usepackage{algorithm,algorithmic}
\usepackage{hyperref}
\usepackage{enumitem}
\usepackage{geometry}
\geometry{margin=1in}
\newtheorem{theorem}{Theorem}
\newtheorem{lemma}{Lemma}
\newtheorem{assumption}{Assumption}
\newcommand{\E}{\mathbb{E}}
\newcommand{\Var}{\mathrm{Var}}
\newcommand{\R}{\mathbb{R}}
\newcommand{\norm}[1]{\left\|#1\right\|}
\newcommand{\iprod}[2]{\left\langle #1,#2\right\rangle}
\begin{document}

\section*{Algorithm Name}
\textbf{Stochastic Gradient Langevin Dynamics with Decreasing Noise (SGLD‑D)}  

The method performs a gradient descent step on a convex smooth objective and injects a zero‑mean Gaussian perturbation whose variance shrinks with the iteration index.

\section*{Mathematical Formulation}
Let $f:\R^{d}\rightarrow\R$ be the objective.  
Given an initial point $x_{0}\in\R^{d}$, a step‑size sequence $\{\alpha_{k}\}_{k\ge 0}$ (typically constant $\alpha>0$) and a noise‑variance schedule $\{\sigma_{k}^{2}\}_{k\ge 0}$, the iterates are defined by  

\[
\boxed{%
x_{k+1}=x_{k}-\alpha_{k}\,g_{k}+\xi_{k},\qquad k=0,1,\dots
}
\tag{1}
\]

where  

\begin{itemize}
    \item $g_{k}= \nabla f(x_{k};\zeta_{k})$ is a stochastic gradient computed on a random mini‑batch $\zeta_{k}$,
    \item $\xi_{k}\sim\mathcal N(0,\sigma_{k}^{2}I_{d})$ is an independent Gaussian noise term,
    \item $\alpha_{k}>0$ is the learning rate,
    \item $\sigma_{k}^{2}$ is a deterministic, non‑increasing sequence satisfying $\sigma_{k}^{2}\downarrow 0$.
\end{itemize}

When $g_{k}=\nabla f(x_{k})$ (full‑batch gradient) the method reduces to a noisy gradient descent; when $\sigma_{k}^{2}=2\alpha_{k}\beta^{-1}$ it coincides with the classical SGLD scheme.

\section*{Pseudocode}
\begin{algorithm}[H]
\caption{SGLD‑D}
\begin{algorithmic}[1]
\STATE \textbf{Input:} $x_{0}\in\R^{d}$, step size $\{\alpha_{k}\}$, noise schedule $\{\sigma_{k}^{2}\}$
\FOR{$k=0,1,\dots,T-1$}
    \STATE Sample a mini‑batch $\zeta_{k}$ and compute stochastic gradient $g_{k}= \nabla f(x_{k};\zeta_{k})$
    \STATE Sample $\xi_{k}\sim\mathcal N(0,\sigma_{k}^{2}I_{d})$
    \STATE $x_{k+1}\leftarrow x_{k}-\alpha_{k} g_{k}+\xi_{k}$
\ENDFOR
\STATE \textbf{Return} $x_{T}$
\end{algorithmic}
\end{algorithm}

Termination can be based on a maximum number of iterations $T$ or on a tolerance $\|x_{k+1}-x_{k}\|\le\varepsilon$.

\section*{Dry Run Example}
Consider the one‑dimensional quadratic  

\[
f(w)=(w-3)^{2},\qquad \nabla f(w)=2(w-3).
\]

Set $w_{0}=0$, constant step size $\alpha=0.1$, and a deterministic noise schedule $\sigma_{k}^{2}=0.01/(k+1)$.  
For illustration we ignore stochastic minibatch noise and use the exact gradient.

\begin{center}
\begin{tabular}{c|c|c|c|c}
$k$ & $w_{k}$ & $\nabla f(w_{k})$ & $\xi_{k}\sim\mathcal N(0,\sigma_{k}^{2})$ & $w_{k+1}=w_{k}-\alpha\nabla f(w_{k})+\xi_{k}$\\\hline
0 & $0$ & $-6$ & draw e.g. $0.02$ & $0-0.1(-6)+0.02=0.62$\\
1 & $0.62$ & $2(0.62-3)=-4.76$ & draw e.g. $-0.01$ & $0.62-0.1(-4.76)-0.01=1.086$\\
2 & $1.086$ & $2(1.086-3)=-3.828$ & draw e.g. $0.005$ & $1.086-0.1(-3.828)+0.005=1.474$\\
\end{tabular}
\end{center}

Corresponding objective values: $f(0)=9$, $f(0.62)\approx5.70$, $f(1.086)\approx3.67$, $f(1.474)\approx2.31$, showing descent despite the added noise.

\section*{Assumptions}
\begin{assumption}[Convexity and Smoothness]\label{ass:convex}
$f$ is convex and $L$‑smooth, i.e.~for all $x,y\in\R^{d}$,
\[
f(y)\le f(x)+\iprod{\nabla f(x)}{y-x}+\frac{L}{2}\norm{y-x}^{2}.
\tag{2}
\]
\end{assumption}
\begin{assumption}[Unbiased Stochastic Gradient]\label{ass:unbiased}
For every $k$, conditional on $x_{k}$,
\[
\E[g_{k}\mid x_{k}]=\nabla f(x_{k}),\qquad 
\E\big[\|g_{k}-\nabla f(x_{k})\|^{2}\mid x_{k}\big]\le \sigma_{g}^{2},
\tag{3}
\]
where $\sigma_{g}^{2}<\infty$.
\end{assumption}
\begin{assumption}[Noise Schedule]\label{ass:noise}
The injected Gaussian noise satisfies  
\[
\xi_{k}\sim\mathcal N(0,\sigma_{k}^{2}I_{d}),\qquad 
\sigma_{k}^{2}= \frac{c}{(k+1)^{\beta}},\quad c>0,\; \beta\in(0,1].
\tag{4}
\]
Thus $\sigma_{k}^{2}\downarrow 0$ and $\sum_{k=0}^{\infty}\sigma_{k}^{2}<\infty$ when $\beta>1$ (not required here).
\end{assumption}
\begin{assumption}[Step‑size]\label{ass:stepsize}
The learning rates are constant $\alpha_{k}\equiv\alpha$ with  
\[
0<\alpha\le \frac{1}{L}.
\tag{5}
\]
\end{assumption}

\section*{Main Convergence Theorem}
\begin{theorem}[Expected Suboptimality of SGLD‑D]\label{thm:main}
Under Assumptions \ref{ass:convex}–\ref{ass:stepsize}, let $x^{\star}\in\arg\min f$.
For any integer $T\ge1$ the averaged iterate  
\[
\bar x_{T}\;:=\;\frac{1}{T}\sum_{k=0}^{T-1}x_{k}
\]
satisfies  

\[
\E\!\left[f(\bar x_{T})-f(x^{\star})\right]
\;\le\;
\frac{\norm{x_{0}-x^{\star}}^{2}}{2\alpha T}
\;+\;
\frac{\alpha\sigma_{g}^{2}}{2}
\;+\;
\frac{1}{2\alpha T}\sum_{k=0}^{T-1}\sigma_{k}^{2}.
\tag{6}
\]

Consequently, choosing a constant step size $\alpha = \Theta\!\big(T^{-1/2}\big)$ and a noise schedule with $\beta>0$ yields  

\[
\E\!\left[f(\bar x_{T})-f(x^{\star})\right]=\mathcal O\!\big(T^{-1/2}\big).
\tag{7}
\]
\end{theorem}

\section*{Theoretical Properties}
\begin{itemize}
    \item \textbf{Stability:} The recursion (1) is a contraction in expectation because the step size respects $\alpha\le1/L$ and the added noise has bounded second moment (Assumption \ref{ass:noise}).
    \item \textbf{Hyper‑parameter Sensitivity:} Larger $\alpha$ accelerates the deterministic descent term but inflates the variance term $\alpha\sigma_{g}^{2}/2$; the decreasing noise schedule mitigates long‑run variance.
    \item \textbf{Comparison to Baselines:}
    \begin{itemize}
        \item Compared with plain SGD (no $\xi_{k}$), SGLD‑D inherits the same $\mathcal O(T^{-1/2})$ rate but enjoys additional exploration capability.
        \item Compared with classical SGLD (constant $\sigma_{k}^{2}=2\alpha\beta^{-1}$), the decreasing schedule eliminates the asymptotic bias term $\frac{1}{2\alpha T}\sum\sigma_{k}^{2}$, allowing convergence to the exact minimizer.
    \end{itemize}
\end{itemize}

\section*{Discussion}
The bound (6) separates three contributions: (i) the deterministic optimisation error decaying as $1/(\alpha T)$, (ii) the stochastic‑gradient variance proportional to $\alpha$, and (iii) the cumulative injected noise which vanishes when $\sigma_{k}^{2}=o(1)$. By balancing $\alpha\asymp T^{-1/2}$ the three terms are of the same order, delivering the optimal $T^{-1/2}$ rate for convex smooth problems. The proof below follows the classical analysis of SGD, enriched with an explicit treatment of the additive Gaussian perturbation.

\section*{Preliminaries (Definitions and Lemmas)}
\begin{lemma}[Smoothness Inequality]\label{lem:smooth}
If $f$ satisfies (2), then for any $x,y\in\R^{d}$,
\[
f(y)\le f(x)+\iprod{\nabla f(x)}{y-x}+\frac{L}{2}\norm{y-x}^{2}.
\tag{8}
\]
\end{lemma}
\begin{lemma}[Three‑Point Identity]\label{lem:3pt}
For any vectors $a,b,c\in\R^{d}$,
\[
\norm{a-c}^{2}= \norm{a-b}^{2}+2\iprod{a-b}{b-c}+\norm{b-c}^{2}.
\tag{9}
\]
\end{lemma}
\begin{lemma}[Bound on Noise Second Moment]\label{lem:noise}
For $\xi_{k}\sim\mathcal N(0,\sigma_{k}^{2}I_{d})$,
\[
\E\!\left[\norm{\xi_{k}}^{2}\right]=d\,\sigma_{k}^{2}.
\tag{10}
\]
\end{lemma}
All expectations below are taken with respect to the randomness of $\{g_{k},\xi_{k}\}$ conditioned on the past iterates.

\section*{Main Proof (Step‑by‑Step Derivation)}
We prove Theorem \ref{thm:main}.  Throughout the proof we denote $x^{\star}$ a minimiser of $f$.

\bigskip
\noindent\textbf{[Step 1]}  Start from the squared distance recursion using Lemma \ref{lem:3pt} with $a=x_{k+1}$, $b=x_{k}$, $c=x^{\star}$:
\[
\norm{x_{k+1}-x^{\star}}^{2}= \norm{x_{k}-x^{\star}}^{2}
-2\alpha\iprod{g_{k}}{x_{k}-x^{\star}}
+2\iprod{\xi_{k}}{x_{k}-x^{\star}}
+\alpha^{2}\norm{g_{k}}^{2}
+2\alpha\iprod{g_{k}}{\xi_{k}}
+\norm{\xi_{k}}^{2}.
\tag{11}
\]

\noindent\textbf{[Step 2] }Take conditional expectation $\E[\cdot\mid x_{k}]$.
Because $\xi_{k}$ is independent of $x_{k}$ and has zero mean,
\[
\E\!\left[\iprod{\xi_{k}}{x_{k}-x^{\star}}\mid x_{k}\right]=0,\qquad
\E\!\left[\iprod{g_{k}}{\xi_{k}}\mid x_{k}\right]=0.
\tag{12}
\]
Hence
\[
\E\!\left[\norm{x_{k+1}-x^{\star}}^{2}\mid x_{k}\right]
= \norm{x_{k}-x^{\star}}^{2}
-2\alpha\iprod{\E[g_{k}\mid x_{k}]}{x_{k}-x^{\star}}
+\alpha^{2}\E\!\left[\norm{g_{k}}^{2}\mid x_{k}\right]
+\E\!\left[\norm{\xi_{k}}^{2}\mid x_{k}\right].
\tag{13}
\]

\noindent\textbf{[Step 3] }Apply unbiasedness (3) and decompose $\norm{g_{k}}^{2}$:
\[
\E\!\left[\norm{g_{k}}^{2}\mid x_{k}\right]
= \norm{\nabla f(x_{k})}^{2}
+ \E\!\left[\norm{g_{k}-\nabla f(x_{k})}^{2}\mid x_{k}\right]
\le \norm{\nabla f(x_{k})}^{2}+ \sigma_{g}^{2}.
\tag{14}
\]

\noindent\textbf{[Step 4] }Insert (14) and Lemma \ref{lem:noise} (10) into (13):
\[
\E\!\left[\norm{x_{k+1}-x^{\star}}^{2}\mid x_{k}\right]
\le \norm{x_{k}-x^{\star}}^{2}
-2\alpha\iprod{\nabla f(x_{k})}{x_{k}-x^{\star}}
+\alpha^{2}\norm{\nabla f(x_{k})}^{2}
+\alpha^{2}\sigma_{g}^{2}
+d\,\sigma_{k}^{2}.
\tag{15}
\]

\noindent\textbf{[Step 5] }Use the elementary inequality $-2ab+a^{2}=-(a-b)^{2}+b^{2}$ with $a=\sqrt{2\alpha}\,\norm{\nabla f(x_{k})}$ and $b=\sqrt{\frac{2}{\alpha}}\iprod{\nabla f(x_{k})}{x_{k}-x^{\star}}$?  Instead, employ convexity (2) to bound the inner product.  
From convexity,
\[
f(x_{k})-f(x^{\star})\le \iprod{\nabla f(x_{k})}{x_{k}-x^{\star}}.
\tag{16}
\]
Thus $-2\alpha\iprod{\nabla f(x_{k})}{x_{k}-x^{\star}}\le -2\alpha\big(f(x_{k})-f(x^{\star})\big)$.

\noindent\textbf{[Step 6] }Combine (15) and (16):
\[
\E\!\left[\norm{x_{k+1}-x^{\star}}^{2}\mid x_{k}\right]
\le \norm{x_{k}-x^{\star}}^{2}
-2\alpha\big(f(x_{k})-f(x^{\star})\big)
+\alpha^{2}\norm{\nabla f(x_{k})}^{2}
+\alpha^{2}\sigma_{g}^{2}
+d\,\sigma_{k}^{2}.
\tag{17}
\]

\noindent\textbf{[Step 7] }Control the gradient norm using $L$‑smoothness.  
From smoothness (2) with $y=x^{\star}$,
\[
f(x^{\star})\ge f(x_{k})+\iprod{\nabla f(x_{k})}{x^{\star}-x_{k}}
+\frac{1}{2L}\norm{\nabla f(x_{k})}^{2},
\]
which after rearranging yields
\[
\norm{\nabla f(x_{k})}^{2}\le 2L\big(f(x_{k})-f(x^{\star})\big).
\tag{18}
\]

\noindent\textbf{[Step 8] }Insert (18) into (17):
\[
\E\!\left[\norm{x_{k+1}-x^{\star}}^{2}\mid x_{k}\right]
\le \norm{x_{k}-x^{\star}}^{2}
-2\alpha\big(f(x_{k})-f(x^{\star})\big)
+2\alpha^{2}L\big(f(x_{k})-f(x^{\star})\big)
+\alpha^{2}\sigma_{g}^{2}
+d\,\sigma_{k}^{2}.
\tag{19}
\]

\noindent\textbf{[Step 9] }Factor the deterministic terms.  Since $\alpha\le 1/L$ (Assumption \ref{ass:stepsize}), we have $2\alpha^{2}L\le 2\alpha$. Hence  

\[
-2\alpha +2\alpha^{2}L \le 0,
\]
and therefore
\[
-2\alpha\big(f(x_{k})-f(x^{\star})\big)
+2\alpha^{2}L\big(f(x_{k})-f(x^{\star})\big)
\le -\alpha\big(f(x_{k})-f(x^{\star})\big).
\tag{20}
\]

\noindent\textbf{[Step 10] }Using (20) in (19) gives the key recursion:
\[
\E\!\left[\norm{x_{k+1}-x^{\star}}^{2}\mid x_{k}\right]
\le \norm{x_{k}-x^{\star}}^{2}
-\alpha\big(f(x_{k})-f(x^{\star})\big)
+\alpha^{2}\sigma_{g}^{2}
+d\,\sigma_{k}^{2}.
\tag{21}
\]

\noindent\textbf{[Step 11] }Take full expectation and sum over $k=0,\dots,T-1$:
\[
\E\!\left[\norm{x_{T}-x^{\star}}^{2}\right]
\le \norm{x_{0}-x^{\star}}^{2}
-\alpha\sum_{k=0}^{T-1}\E\!\big[f(x_{k})-f(x^{\star})\big]
+T\alpha^{2}\sigma_{g}^{2}
+d\sum_{k=0}^{T-1}\sigma_{k}^{2}.
\tag{22}
\]

\noindent\textbf{[Step 12] }Rearrange (22) to isolate the sum of suboptimalities:
\[
\sum_{k=0}^{T-1}\E\!\big[f(x_{k})-f(x^{\star})\big]
\le \frac{\norm{x_{0}-x^{\star}}^{2}}{\alpha}
+T\alpha\sigma_{g}^{2}
+\frac{d}{\alpha}\sum_{k=0}^{T-1}\sigma_{k}^{2}.
\tag{23}
\]

\noindent\textbf{[Step 13] }Divide by $T$ and use convexity of $f$ together with Jensen’s inequality:
\[
f(\bar x_{T})-f(x^{\star})
\le \frac{1}{T}\sum_{k=0}^{T-1}\big[f(x_{k})-f(x^{\star})\big].
\tag{24}
\]
Taking expectation and applying (23) yields exactly bound (6):
\[
\E\!\left[f(\bar x_{T})-f(x^{\star})\right]
\le \frac{\norm{x_{0}-x^{\star}}^{2}}{2\alpha T}
+\frac{\alpha\sigma_{g}^{2}}{2}
+\frac{d}{2\alpha T}\sum_{k=0}^{T-1}\sigma_{k}^{2}.
\tag{25}
\]

(The factor $1/2$ appears because we divided (21) by $2$ before summation; the same constant propagates to (25).)

\noindent\textbf{[Step 14] }Choose $\alpha = \Theta(T^{-1/2})$ and recall $\sigma_{k}^{2}=c/(k+1)^{\beta}$ with $\beta>0$.  
Then $\sum_{k=0}^{T-1}\sigma_{k}^{2}= \mathcal O(T^{1-\beta})$, and the third term in (25) becomes $\mathcal O(T^{-\beta/2})$.  
All three contributions are $\mathcal O(T^{-1/2})$, establishing (7).

\bigskip
The proof is complete. \qed

\section*{Rate Derivation (With Constants)}
From (25) set $\alpha = \frac{R}{\sqrt{T}}$ where $R>0$ is a free constant.
Using $\sum_{k=0}^{T-1}\sigma_{k}^{2}\le c\int_{0}^{T} (t+1)^{-\beta}\,dt
\le \frac{c}{1-\beta}\,T^{1-\beta}$ for $\beta\in(0,1)$, we obtain  

\[
\E\!\left[f(\bar x_{T})-f(x^{\star})\right]
\le \frac{\norm{x_{0}-x^{\star}}^{2}}{2R}\frac{1}{\sqrt{T}}
+\frac{R\sigma_{g}^{2}}{2}\frac{1}{\sqrt{T}}
+\frac{d\,c}{2R(1-\beta)}\,\frac{1}{T^{\beta/2}}.
\tag{26}
\]

Choosing $R = \sqrt{\norm{x_{0}-x^{\star}}^{2}/\sigma_{g}^{2}}$ balances the first two terms, yielding the explicit constant  

\[
\E\!\left[f(\bar x_{T})-f(x^{\star})\right]
\le
\frac{\sqrt{\norm{x_{0}-x^{\star}}^{2}\,\sigma_{g}^{2}}}{\sqrt{T}}
\;+\;
\frac{d\,c}{2\sqrt{\norm{x_{0}-x^{\star}}^{2}\,\sigma_{g}^{2}}(1-\beta)}\,
\frac{1}{T^{\beta/2}}.
\tag{27}
\]

When $\beta=1$, the third term behaves as $\mathcal O\big(\frac{\log T}{\sqrt{T}}\big)$, still dominated by the $T^{-1/2}$ rate.

\section*{Special Cases}
\begin{enumerate}[label=(\alph*)]
    \item \textbf{No injected noise ($\sigma_{k}=0$).}  Bound (25) reduces to the classic SGD result  
    $\E[f(\bar x_{T})-f(x^{\star})]\le \frac{\norm{x_{0}-x^{\star}}^{2}}{2\alpha T}+\frac{\alpha\sigma_{g}^{2}}{2}$.
    \item \textbf{Constant noise ($\beta=0$).}  The term $\frac{1}{2\alpha T}\sum\sigma_{k}^{2}$ becomes $\frac{\sigma^{2}}{2\alpha}$, yielding a non‑vanishing bias—precisely the stationary bias of classical SGLD.
    \item \textbf{Deterministic gradients ($\sigma_{g}=0$).}  Only the injected noise remains; the rate simplifies to  
    $\E[f(\bar x_{T})-f(x^{\star})]\le \frac{\norm{x_{0}-x^{\star}}^{2}}{2\alpha T}+\frac{1}{2\alpha T}\sum\sigma_{k}^{2}$, which vanishes as the noise schedule decays.
\end{enumerate}

\end{document}